\documentclass[12pt]{article}
\usepackage[margin=1in]{geometry}
\usepackage{amsmath,amssymb}
\usepackage{enumitem}
\usepackage{array}
\usepackage{booktabs}

\begin{document}

\begin{center}
    {\Large \textbf{SOLUTIONS TO THE PRACTICE TEST FOR TEST 1 (STAT-461/561)}}
\end{center}

\begin{enumerate}[label=(\arabic*), leftmargin=*]

\item
\begin{enumerate}[label=(\alph*)]
    \item F
    \item F; the order is wrong.
    \item F; simple random sampling gets rid of selection bias but not other types of bias such as non-response bias.
    \item F
    \item F
    \item F
    \item T
    \item F
    \item T
    \item F
    \item T
    \item T
    \item F
    \item F
\end{enumerate}

\item
\begin{enumerate}[label=(\alph*)]
    \item Continuous ratio-type quantitative (or numerical).
    \item Interval-type quantitative.
    \item Ordinal qualitative (or categorical).
    \item Ordinal qualitative (or categorical).
    \item Nominal qualitative (or categorical).
    \item Discrete ratio-type quantitative.
    \item Interval-type quantitative.
\end{enumerate}

\item
\begin{enumerate}[label=(\alph*)]
    \item Convenience sampling.
    \item Systematic sampling.
    \item Cluster sampling.
    \item Stratified random sampling.
\end{enumerate}

\item
\[
\bar{x}=2.75.
\]
The median is
\[
\frac{0+1}{2}=0.5.
\]
The lower quartile is
\[
Q_1=\frac{0+0}{2}=0,
\]
and the upper quartile is
\[
Q_3=\frac{5+6}{2}=5.5.
\]
Therefore,
\[
\operatorname{IQR}=Q_3-Q_1=5.5-0=5.5.
\]
The histogram will vary from person to person.

\item
\begin{enumerate}[label=(\alph*)]
    \item Stem-and-leaf plot:
\begin{center}
\begin{tabular}{r|l}
10 & 9\\
11 & 8\quad 9\\
12 & 0\quad 1\quad 2\quad 2\quad 4\quad 6\quad 7\\
13 & 2\quad 4\quad 5\\
14 & 1\\
15 & \\
16 & \\
17 & \\
18 & \\
19 & \\
20 & \\
21 & \\
22 & \\
23 & 1
\end{tabular}
\end{center}

    \item The five-point summary is
    \[
    \min=10.9,\quad Q_1=12.0,\quad \text{Median}=12.4,\quad Q_3=13.4,\quad \max=23.1.
    \]

    \item For the modified boxplot, the value $23.1$ is an upper outlier. The box is based on $Q_1=12.0$, median $=12.4$, and $Q_3=13.4$, with the upper non-outlier whisker at $14.1$.
\end{enumerate}

\item
If the data-value $23.1$ becomes $13.1$, the mean, range, and standard deviation will change significantly, whereas the median, IQR, and MAD will not. The median, IQR, and MAD are resistant to changes in the extremes of the dataset.

For the original dataset,
\[
\bar{x}=13.207,\qquad s=2.851,\qquad \text{Range}=23.1-10.9=12.2.
\]
For the corrected dataset,
\[
\bar{x}=12.54,\qquad s=0.813,\qquad \text{Range}=14.1-10.9=3.2.
\]

\item
The standard deviation of the dataset in problem (5) is
\[
s=2.851.
\]
The MAD is the median of the absolute deviations from the overall median $12.4$:
\[
|12.4-12.4|,\ |12.2-12.4|,\ |13.4-12.4|,\ |10.9-12.4|,\ldots,\ |12.6-12.4|,\ |11.9-12.4|,\ |23.1-12.4|.
\]
Thus the MAD is the median of
\[
0,\ 0.2,\ 1.0,\ 1.5,\ldots,\ 0.2,\ 0.5,\ 10.7,
\]
so
\[
\operatorname{MAD}=0.5.
\]

\item
The sample correlation is
\[
r=0.941.
\]
The equation of the least-squares line is
\[
Y=0.561+0.5919X.
\]
The coefficient of determination is
\[
r^2=0.8836=88.36\%.
\]
Since the best-fitting line explains $88.36\%$ of the overall variation in $Y$, the line fits the dataset well.

\end{enumerate}

\end{document}
